% !TeX root = ../thuthesis-example.tex

\chapter{基于协同引导的多模态语义ID构建方法}

本章首先介绍生成式推荐的两阶段框架，阐述语义ID构建与生成式推荐的关系；然后详细介绍本文提出的COMET（Collaborative-guided Multimodal Encoding Tokenizer）方法，包括问题定义、整体框架、多模态嵌入提取、多模态融合模块、残差向量量化、损失函数设计以及Sinkhorn碰撞优化策略。

\section{生成式推荐的两阶段框架}

生成式推荐将传统的序列推荐任务重新表述为序列到序列的生成问题，其核心思想是将物品用离散的语义ID表示，然后使用生成模型预测用户下一个交互物品的语义ID。整个框架分为两个阶段：\textbf{语义ID构建}和\textbf{生成式推荐}。

\subsection{第一阶段：语义ID构建}

语义ID构建的目标是为每个物品学习一个多层离散编码，该编码应满足：（1）语义相似的物品具有相近的编码；（2）不同物品尽量具有不同的编码；（3）编码层级从粗到细，捕捉从类别到细节的层次化语义信息。

形式化地，设物品集合为$\mathcal{I} = \{i_1, i_2, \ldots, i_N\}$，语义ID构建过程可表示为映射函数$f: \mathcal{I} \to \mathcal{C}^L$，其中$\mathcal{C} = \{0, 1, \ldots, K-1\}$为单层码本空间，$L$为码本层数。对于物品$i$，其语义ID为：
\begin{equation}
  \text{SID}(i) = (c_1, c_2, \ldots, c_L), \quad c_l \in \mathcal{C}.
\end{equation}

现有方法主要通过向量量化技术实现语义ID构建。典型流程为：首先使用预训练的嵌入模型（如文本编码器、多模态编码器）提取物品的语义表示，然后通过残差向量量化（RQ-VAE）或层次化聚类方法将连续的嵌入向量映射为多层离散编码。本文的COMET方法属于这一阶段，聚焦于通过协同引导的多模态融合提升语义ID的质量。

\subsection{第二阶段：生成式推荐}

在语义ID构建完成后，生成式推荐将序列推荐转化为条件生成任务。给定用户的交互历史序列$\mathcal{S}_u = [s_1, s_2, \ldots, s_T]$，将每个物品替换为其语义ID，得到输入序列：
\begin{equation}
  \mathbf{X} = [\text{SID}(s_1); \text{SID}(s_2); \ldots; \text{SID}(s_T)].
\end{equation}

生成模型以$\mathbf{X}$为输入，自回归地生成目标物品$s_{T+1}$的语义ID。由于语义ID由$L$层编码组成，生成过程可以分解为逐层预测：
\begin{equation}
  P(s_{T+1} | \mathcal{S}_u) = \prod_{l=1}^{L} P(c_l | \mathbf{X}, c_1, \ldots, c_{l-1}).
\end{equation}

\textbf{两阶段框架的意义}。将语义ID构建与生成式推荐分离为两个独立阶段具有重要意义：第一，语义ID作为中间表示，解耦了物品内容理解与用户行为建模，使得第一阶段可以专注于挖掘物品的多模态语义特征，第二阶段可以专注于捕捉用户的序列行为模式；第二，语义ID的离散性使得生成模型能够利用成熟的序列生成技术（如Transformer、束搜索等），并天然支持物品的层次化检索；第三，高质量的语义ID能够为下游生成模型提供更强的归纳偏置——语义相似的物品拥有相近的编码，使得模型更容易学习到用户偏好在语义空间中的迁移规律，从而提升推荐的准确性和泛化能力。

然而，语义ID的质量直接决定了整个框架的性能上限。如果语义ID不能有效编码物品的多模态特征和协同信号，即使第二阶段的生成模型再先进，也难以学习到有效的推荐模式。因此，本文聚焦于第一阶段的语义ID构建，提出COMET方法，通过协同引导的多模态融合机制提升语义ID的表示能力，为下游生成式推荐奠定坚实基础。

\section{COMET方法概述}

\subsection{问题定义}

本文聚焦于语义ID构建任务（定义见3.1.1节）。设物品集合为$\mathcal{I} = \{i_1, i_2, \ldots, i_N\}$，用户集合为$\mathcal{U} = \{u_1, u_2, \ldots, u_M\}$。

COMET的目标是学习一个编码器$f_{\theta}: \mathcal{I} \to \mathcal{C}^L$，通过融合物品的文本、图像和协同过滤信息构建高质量的多层离散语义ID，使其同时反映物品的内容语义和协同过滤信号。在本文的实验设置中，码本层数$L=4$，每层码本大小$K=256$。

\subsection{整体框架}

COMET的整体框架如图~\ref{fig:framework}所示，主要包括以下几个部分：

\begin{enumerate}
  \item \textbf{多模态嵌入提取}：分别使用Qwen3-Embedding-8B、CLIP ViT-L/14和SASRec提取物品的文本嵌入（4096维）、图像嵌入（768维）和协同过滤嵌入（128维）。
  \item \textbf{多模态融合模块}：通过协同引导交叉注意力机制将三种模态嵌入融合为64维隐向量。
  \item \textbf{残差向量量化}：将隐向量通过4层残差量化生成离散的4层语义ID。
  \item \textbf{下游生成推荐}：将构建的语义ID输入TIGER或OneRec进行生成式推荐。
\end{enumerate}

\begin{figure}[htbp]
  \centering
  \includegraphics[width=\textwidth]{framework-comet.png}
  \caption{COMET整体框架}
  \label{fig:framework}
\end{figure}

\subsection{多模态嵌入提取}

语义ID构建的质量取决于输入特征的表达能力。单一模态的信息往往不足以完整刻画物品的特性：仅依赖文本无法精确描述视觉外观（如颜色、质地），而仅依赖图像则丧失了品牌、价格等结构化属性。更重要的是，这些内容模态都缺乏用户行为维度的信息——它们无法反映"哪些物品经常被一起购买"这类协同过滤信号。因此，COMET的输入由三种模态的物品嵌入组成，分别编码物品的文本语义、视觉外观和协同过滤信号，为后续的融合模块提供互补的信息源。

\subsubsection{文本嵌入}

文本嵌入是语义ID的核心信息来源，它承载了物品最显式、最结构化的语义信息。为充分利用电商场景中的多字段文本数据，本文为每个物品构建结构化的文本描述，包含标题、品牌、价格、商品描述和类别信息，格式如下：
\begin{verbatim}
Title: <标题>
Brand: <品牌>
Price: $<价格>
Description: <商品描述>
Categories: <类别1> > <类别2> > <类别3>
\end{verbatim}

使用Qwen3-Embedding-8B\cite{qwen3embedding}作为文本编码器。该模型是一个80亿参数的文本嵌入模型，通过last-token池化提取4096维的密集表示。在推理时，为输入文本添加指令前缀``Represent the item description for recommendation:''以引导模型关注推荐相关的语义信息。

\subsubsection{图像嵌入}

图像嵌入补充了物品的视觉外观信息。在电商场景中，商品图片直观展示了颜色、形状、材质等文本难以精确描述的属性，对用户的购买决策具有重要影响。例如，在化妆品推荐中，两款唇膏可能具有几乎相同的文本描述（品牌、系列、价格），但颜色差异显著（橙色 vs. 棕色），用户的选择往往取决于视觉外观。仅依赖文本嵌入会导致这类物品被分配相同的语义ID，无法满足用户对特定视觉属性的偏好。

本文使用CLIP ViT-L/14\cite{radford2021learning}作为图像编码器。CLIP通过在4亿图文对上进行对比学习，其视觉编码器能够提取具有语义理解能力的768维图像表示。对于少量缺失图像的物品（Amazon Beauty数据集中约8个物品），使用零向量作为占位，并在多模态融合模块中通过可学习填充向量处理。

\subsubsection{协同过滤嵌入}

协同过滤嵌入编码了物品在用户行为空间中的表示，反映了"购买了A的用户也倾向于购买B"这类协同信号。与文本和图像这两种内容模态不同，协同过滤嵌入完全来自用户的实际交互数据，它捕捉的是物品在推荐场景下的真实关联模式——这些模式可能源于功能互补（如洗面奶和爽肤水）、使用场景相似（如运动鞋和运动服）或用户群体重叠，而不一定能从内容特征中直接推断。引入协同过滤信号的意义在于：使语义ID的构建不仅基于物品"是什么"（内容），还考虑物品"如何被使用"（行为），从而更好地服务于推荐任务。

本文使用预训练的SASRec\cite{kang2018self}模型提取协同过滤嵌入。SASRec在用户交互序列上通过交叉熵损失训练，其物品嵌入表$\mathbf{E} \in \mathbb{R}^{N \times d_{\text{cf}}}$中的每一行编码了对应物品的协同过滤信息。本文使用的SASRec模型的隐藏维度$d_{\text{cf}}=128$。

\subsection{基于协同引导的多模态融合模块}

拥有三种模态的嵌入后，关键问题是如何有效融合它们。简单的拼接或加权平均无法建模模态间的复杂交互关系，也无法根据物品特性自适应地调整各模态的重要性。COMET的核心创新是协同引导交叉注意力融合模块，其设计理念基于以下观察：协同过滤信号反映了用户群体的真实行为偏好，它本质上定义了"什么样的物品在推荐场景下是相似的"。因此，以协同过滤嵌入作为查询（Query），让其从内容特征（文本和图像）中自适应地提取推荐相关信息，能够实现架构级别的协同引导——内容模态通过注意力权重被动态地重新组织，以最大化与协同信号的对齐。

\subsubsection{特征对齐投影器}

三种模态的原始嵌入维度差异巨大（文本4096维、图像768维、协同过滤128维），且来自不同的预训练模型，其表示空间的尺度和语义结构各不相同。直接将它们输入交叉注意力会导致维度不匹配和语义空间不对齐的问题。因此，首先需要通过特征对齐投影器将三种模态嵌入投影到统一的注意力操作维度$d_{\text{model}}$（本文设为256）。

\textbf{文本多token投影}。原始文本嵌入$\mathbf{e}_{\text{text}} \in \mathbb{R}^{4096}$维度远高于其他模态，单一线性投影（$4096 \to 256$）会导致严重的信息压缩损失——压缩比高达16倍。为缓解这一问题，COMET将4096维文本嵌入均匀分割为4个1024维的片段，每个片段使用独立的线性投影映射到$d_{\text{model}}$维：
\begin{equation}
  \mathbf{e}_{\text{text}} = [\mathbf{e}_{\text{text}}^{(1)}; \mathbf{e}_{\text{text}}^{(2)}; \mathbf{e}_{\text{text}}^{(3)}; \mathbf{e}_{\text{text}}^{(4)}], \quad \mathbf{e}_{\text{text}}^{(j)} \in \mathbb{R}^{1024},
\end{equation}
\begin{equation}
  \mathbf{h}_{\text{text}}^{(j)} = \mathbf{W}_{\text{text}}^{(j)} \mathbf{e}_{\text{text}}^{(j)} + \mathbf{b}_{\text{text}}^{(j)}, \quad \mathbf{W}_{\text{text}}^{(j)} \in \mathbb{R}^{d_{\text{model}} \times 1024}, \quad j = 1, 2, 3, 4.
\end{equation}
这一设计将原始的1个文本KV token扩展为4个token，每个token聚焦于文本嵌入的不同语义子空间，使交叉注意力能够更细粒度地选择文本信息。

图像嵌入和协同过滤嵌入的投影方式不变：
\begin{align}
  \mathbf{h}_{\text{image}} &= \mathbf{W}_{\text{image}} \mathbf{e}_{\text{image}} + \mathbf{b}_{\text{image}}, \quad \mathbf{W}_{\text{image}} \in \mathbb{R}^{d_{\text{model}} \times 768}, \\
  \mathbf{h}_{\text{cf}} &= \mathbf{W}_{\text{cf}} \mathbf{e}_{\text{cf}} + \mathbf{b}_{\text{cf}}, \quad \mathbf{W}_{\text{cf}} \in \mathbb{R}^{d_{\text{model}} \times 128}.
\end{align}

对于缺失图像的物品，使用可学习的填充向量$\mathbf{p} \in \mathbb{R}^{768}$替代零向量：
\begin{equation}
  \mathbf{e}_{\text{image}}' = \begin{cases} \mathbf{e}_{\text{image}}, & \text{若图像存在,} \\ \mathbf{p}, & \text{否则.} \end{cases}
\end{equation}

\subsubsection{协同引导交叉注意力网络}

模态对齐后，需要决定如何组织融合过程。COMET的关键设计是将协同过滤嵌入作为注意力查询（Query），文本和图像嵌入作为键值对（Key-Value），通过多层Transformer解码器块实现深度融合。这一设计的动机在于：Query决定了"寻找什么信息"，而Key-Value提供"可选择的信息源"。将协同过滤信号作为Query，意味着让用户行为模式主导信息提取过程——对于每个物品，根据其在协同过滤空间中的表示，自适应地从内容特征中选择最相关的语义成分。例如，对于经常与护肤品一起购买的化妆品，模型可能更关注其功效描述（文本）；而对于主要依赖视觉吸引力的装饰品，模型可能更关注其外观（图像）。这种架构级的协同引导比LETTER的辅助损失方式更直接——融合过程本身就在协同信号的指导下进行，无需额外的对比损失函数。

\textbf{查询（Query）}：投影后的CF嵌入$\mathbf{Q} = \mathbf{h}_{\text{cf}} \in \mathbb{R}^{1 \times d_{\text{model}}}$。

\textbf{键值对（Key-Value）}：4个文本token和1个图像token拼接构成5个token的序列：
\begin{equation}
  \mathbf{K} = \mathbf{V} = [\mathbf{h}_{\text{text}}^{(1)}; \mathbf{h}_{\text{text}}^{(2)}; \mathbf{h}_{\text{text}}^{(3)}; \mathbf{h}_{\text{text}}^{(4)}; \mathbf{h}_{\text{image}}] \in \mathbb{R}^{5 \times d_{\text{model}}}.
\end{equation}

\textbf{多层交叉注意力}。COMET使用2层堆叠的Transformer解码器块，每层包含交叉注意力和前馈网络（FFN）。设第$i$层输入为$\mathbf{q}^{(i)}$（$\mathbf{q}^{(0)} = \mathbf{Q}$），每层的计算如下：
\begin{align}
  \mathbf{q}'^{(i)} &= \mathbf{q}^{(i)} + \text{CrossAttn}^{(i)}\bigl(\text{LN}_1^{(i)}(\mathbf{q}^{(i)}),\; \text{LN}_2^{(i)}(\mathbf{KV})\bigr), \label{eq:preln-attn} \\
  \mathbf{q}^{(i+1)} &= \mathbf{q}'^{(i)} + \text{FFN}^{(i)}\bigl(\text{LN}_3^{(i)}(\mathbf{q}'^{(i)})\bigr), \label{eq:preln-ffn}
\end{align}
其中$\text{LN}$为层归一化，$\text{FFN}$为两层线性变换（维度$d_{\text{model}} \to 2 d_{\text{model}} \to d_{\text{model}}$）加GELU激活。交叉注意力使用$n_{\text{heads}}=4$个头，$d_{\text{head}}=64$。

与单层注意力相比，2层堆叠结构使模型具有更强的模态交互建模能力。第一层可以学习初步的模态选择和信息提取，第二层在此基础上进一步精细融合。每层的残差连接保证了CF信号的直接传递。

最终融合输出$\mathbf{f} = \mathbf{q}^{(2)}$经过最终的层归一化后输入码本特征映射器。

\subsubsection{码本特征映射器}

归一化后的特征通过两层感知机映射到码本维度$d_e$：
\begin{equation}
  \mathbf{z} = \mathbf{W}_2 \cdot \text{GELU}(\mathbf{W}_1 \mathbf{f} + \mathbf{b}_1) + \mathbf{b}_2,
\end{equation}
其中$\mathbf{W}_1 \in \mathbb{R}^{256 \times 128}$，$\mathbf{W}_2 \in \mathbb{R}^{128 \times 64}$，最终输出$\mathbf{z} \in \mathbb{R}^{64}$。

\subsection{残差向量量化}

融合后的隐向量是连续的高维表示，而生成式推荐需要离散的符号序列作为输入。向量量化是连接连续表示空间与离散符号空间的桥梁。COMET采用4层残差向量量化将融合后的64维隐向量$\mathbf{z}$映射为离散的语义ID。

残差量化的优势在于其层次化特性：第一层量化捕捉物品的粗粒度语义（如类别），后续各层对前一层的残差进行量化，逐步细化编码，实现从粗到细的层次化表示。量化过程与第2章介绍的RQ-VAE相同：第一层对隐向量进行量化，后续各层对前一层的残差进行量化，最终得到4层编码$(c_1, c_2, c_3, c_4)$。这种层次化结构使得生成模型能够逐层生成语义ID，先确定大类别，再逐步细化到具体物品，与人类认知过程相符。

由于向量量化操作不可微，训练时采用直通估计器技术：前向传播使用离散的量化向量$\hat{\mathbf{z}}$，反向传播时梯度直接穿过量化操作传递给编码器。这使得编码器能够通过重建损失进行端到端优化。

为保证码本的有效利用，COMET采用了以下技术：训练开始前使用KMeans对各层码本进行初始化；训练过程中码本通过指数移动平均方式更新以保持稳定性，并对长期未被使用的码向量进行重启。

\subsection{损失函数设计}

COMET的训练损失由重建损失$\mathcal{L}_{\text{rec}}$和量化损失$\mathcal{L}_{\text{vq}}$组成：
\begin{equation}
  \mathcal{L} = \mathcal{L}_{\text{rec}} + \lambda \mathcal{L}_{\text{vq}},
  \label{eq:loss}
\end{equation}
其中$\lambda = 1.0$为量化损失权重。

重建损失采用多目标重建策略，同时重建文本、图像和协同过滤三种模态的嵌入。COMET使用三个独立的解码器分支：

\begin{itemize}
  \item \textbf{文本解码器}：5层MLP（$64 \to 256 \to 512 \to 1024 \to 2048 \to 4096$），重建原始文本嵌入；
  \item \textbf{图像解码器}：2层MLP（$64 \to 256 \to 768$），重建原始图像嵌入；
  \item \textbf{CF解码器}：2层MLP（$64 \to 128 \to 128$），重建原始协同过滤嵌入。
\end{itemize}
中间层使用ReLU激活。设$\text{Dec}_{\text{text}}(\cdot)$、$\text{Dec}_{\text{image}}(\cdot)$和$\text{Dec}_{\text{cf}}(\cdot)$分别表示文本、图像和协同过滤解码器函数。

重建损失为文本、图像和协同过滤三个分支的加权均方误差之和：
\begin{equation}
  \mathcal{L}_{\text{rec}} = w_{\text{text}} \|\text{Dec}_{\text{text}}(\hat{\mathbf{z}}) - \mathbf{e}_{\text{text}}\|^2 + w_{\text{image}} \|\text{Dec}_{\text{image}}(\hat{\mathbf{z}}) - \mathbf{e}_{\text{image}}\|^2 + w_{\text{cf}} \|\text{Dec}_{\text{cf}}(\hat{\mathbf{z}}) - \mathbf{e}_{\text{cf}}\|^2,
\end{equation}
其中$\hat{\mathbf{z}}$为经直通估计器处理后的量化表示，$w_{\text{text}} = 1.0$，$w_{\text{image}} = 0.1$，$w_{\text{cf}} = 0.1$。对于缺失图像的物品，图像重建损失项被掩码为零，不参与梯度计算。

多目标重建的设计动机如下：文本嵌入作为维度最高、信息量最丰富的模态，赋予最大的重建权重（1.0），确保隐空间以文本语义为主干。图像和CF的重建权重较低（0.1），起到正则化作用——迫使隐向量$\mathbf{z}$不仅编码文本信息，还需保留一定的视觉和协同过滤信号，从而实现更均衡的多模态融合。相比仅重建文本嵌入，多目标重建在64维隐空间中能更有效地组织多模态信息。

量化损失用于优化残差向量量化过程：
\begin{equation}
  \mathcal{L}_{\text{vq}} = \frac{1}{L} \sum_{l=1}^{L} \|\mathbf{z}_q^{(l)} - \mathbf{z}^{(l)}\|^2,
\end{equation}
其中$\mathbf{z}^{(l)}$为第$l$层的输入向量，$\mathbf{z}_q^{(l)}$为量化后的向量。该损失促使输入向量靠近其被分配到的码向量，从而稳定量化过程并提高码本的利用率。

在训练过程中，使用加权重建损失$\mathcal{L}_{\text{rec}}$作为模型选择标准，保存该损失最低的模型检查点。这是因为量化损失的波动较大，直接使用总训练损失可能导致选择到重建质量次优的模型。

\subsection{Sinkhorn碰撞优化}

训练完成后，理想情况下每个物品应拥有唯一的语义ID。然而，由于码本容量有限（$K^L = 256^4 \approx 4.3$亿个可能的编码，而实际物品数量远小于此），且量化过程基于距离度量，语义极度相似的物品可能被分配到完全相同的编码，形成碰撞（collision）。碰撞的危害在于：生成模型输出的是语义ID而非物品ID，当多个物品共享同一语义ID时，下游推荐系统通常采用启发式规则（如选择交互频次最高者）将ID映射回物品，这意味着碰撞组中除频次最高者外的所有物品永远不会被推荐，直接导致不可恢复的精度损失。因此，降低碰撞率是提升语义ID质量的关键目标。

训练完成后，对所有物品进行编码，统计语义ID的碰撞情况（多个物品共享相同的语义ID）。对碰撞组进行迭代优化：

在最后两层量化器上启用Sinkhorn-Knopp算法\cite{cuturi2013sinkhorn}，将码分配问题转化为最优传输问题。其核心思想是：在保持物品与码向量距离关系的前提下，尽可能均匀地将物品分配到不同的码上，避免多个物品挤占同一个码。对归一化的距离矩阵$\mathbf{D}$进行迭代行列归一化：
\begin{equation}
  \mathbf{Q}^{(0)} = \exp(-\mathbf{D} / \epsilon), \quad \mathbf{Q}^{(t)} = \text{RowNorm}(\text{ColNorm}(\mathbf{Q}^{(t-1)})),
\end{equation}
迭代50次后，$\mathbf{Q}$趋近于均匀分配矩阵，碰撞组中的物品被分配到不同的码上。

与之前仅在最后一层使用统一温度参数的方案不同，COMET采用\textbf{渐进式温度策略}：第3层$\epsilon_3 = 0.005$，第4层$\epsilon_4 = 0.01$。较浅层使用较小的$\epsilon$以保持编码的精确性，较深层使用较大的$\epsilon$以允许更灵活的码分配。在最后两层而非仅最后一层启用Sinkhorn优化，能够在更大的编码空间中进行碰撞消解，显著提高碰撞消除的有效性。


\section{本章小结}

本章首先介绍了生成式推荐的两阶段框架，阐明了语义ID构建与生成式推荐的关系及其重要意义。两阶段分离的设计使得语义ID构建可以专注于挖掘物品的多模态语义特征，而生成模型可以专注于捕捉用户的序列行为模式。语义ID的质量直接决定了整个生成式推荐框架的性能上限。

在此基础上，本章详细介绍了COMET方法的设计细节。COMET的核心创新在于协同引导交叉注意力融合机制，将协同过滤引导从辅助损失的间接方式提升为架构级别的直接融合，同时引入文本多token投影和图像模态信息，构建了5个KV token的丰富键值序列。在训练策略上，COMET采用多目标重建损失（文本+图像+CF）促进多模态信息在隐空间中的均衡编码，使用余弦退火学习率调度和训练时Sinkhorn均匀化进一步提升码本质量。与LETTER相比，COMET通过架构设计实现CF引导机制的融合，不依赖额外的对比损失函数。
